%%%%%%%%%%%%%%%%%%%%%%%%%%%%%%%%%%%%%%%%%%%%%%%%%%%%%%%%%%%%%%%%%%%%%%%%%%%%%%%%
%                                                                              %
% BlackArch Linux Guide                                                        %
%                                                                              %
%%%%%%%%%%%%%%%%%%%%%%%%%%%%%%%%%%%%%%%%%%%%%%%%%%%%%%%%%%%%%%%%%%%%%%%%%%%%%%%%

\documentclass[a4paper, oneside, 11pt]{book}

%%% INCLUDES %%%
\renewcommand{\familydefault}{\sfdefault}

\usepackage{array}
\usepackage{color}
\usepackage{enumerate}
\usepackage{fancyhdr}
\usepackage{fancyvrb}
\usepackage{geometry}
\usepackage{graphicx}
\usepackage{html}
\usepackage{hyperref}
\usepackage{ifpdf}
\usepackage{listings}
\usepackage{pstricks}
\usepackage{supertabular}
\usepackage{tocloft}
\usepackage[utf8]{inputenc}

%%% LOCALIZATION %%%
\renewcommand{\contentsname}{İçindekiler}
\renewcommand{\chaptername}{Bölüm}

%%% LAYOUT %%%
\setlength{\parindent}{0em}
\setlength{\parskip}{1.5ex plus0.5ex minus0.5ex}
\geometry{left=2.5cm, textwidth=16cm, top=3cm, textheight=25cm, bottom=3cm}
\widowpenalty=2000
\clubpenalty=1000
\frenchspacing
\sloppy
\pagecolor[HTML]{FFFFFF}
\color[HTML]{333333}
\setcounter{tocdepth}{10}
\setcounter{secnumdepth}{10}

\hypersetup{
  pdftitle={BlackArch Linux, BlackArch Linux Rehberi},
  pdfsubject={BlackArch Linux, BlackArch Linux Rehberi},
  pdfauthor={BlackArch Linux, BlackArch Linux},
  pdfkeywords={BlackArch Linux, Sızma Testi, Güvenlik, Hacking, Linux},
  pdfcenterwindow=true,
  colorlinks=true,
  breaklinks=true,
  linkcolor=red,
  menucolor=red,
  urlcolor=red
}

% syntax highlighting
% all options should be set here document wide
\lstset{
backgroundcolor=\color[HTML]{EEEEEE},
frame=single,
basicstyle=\footnotesize\ttfamily,
float,
deletekeywords={return},
otherkeywords={mkdir, curl, sudo, sha1sum, grep, cut, sort, wget, makepkg,
pacman, blackman},
keywordstyle=\color{orange},
commentstyle=\color{blue},
stringstyle=\color{red},
language=bash,
showspaces=false,
showtabs=false,
tabsize=2
}

%%% HEADER / FOOTER %%%
\setlength{\headheight}{33pt}
\setlength{\headsep}{33pt}
\lhead{{\includegraphics[width=1cm,height=1cm]{images/logo.png}}}
\rhead{The BlackArch Linux Rehberi}

%%% CUSTOM MACROS %%%
% for html links
\ifpdf\else
\def\href#1#2{\htmladdnormallink{#2}{#1}}
\fi

%------------------%
%  TITLE PAGE      %
%------------------%
\begin{document}
\pagestyle{empty}
\begin{center}
\begin{figure}[htbp]
\centering
\vspace{0.5cm}
\includegraphics[width=8cm]{images/logo.png}
\label{fig:logo}
\end{figure}
\vspace{0.5cm}
\Huge{\textbf{BlackArch Linux Rehberi}}\\
\vspace{1cm}
\Large{\color{red}https://www.blackarch.org/}\\
\vspace{0.5cm}
\end{center}
\newpage
\tableofcontents
\newpage
\pagestyle{fancy}

%------------------%
%  Chapter 1       %
%------------------%

\chapter{Giriş}

\section{Ön Bakış}
Bu rehber 4 ana parçaya bölünmüştür:
\begin{itemize}
\item Giriş - Projeye genel bir bakış, giriş ve çeşitli yararlı ilk bilgiler
\item Kullanıcı Rehberi - Normal bir kullanıcının BlackArch'ı verimli kullanması için gerekli bilgiler
\item Geliştirici Rehberi - BlackArch'ın geliştirilmesine yardımcı veya destek olmak isteyenler için gerekli bilgiler
\item Araç Rehberi - Kurulu araçların kullanım yöntemleri ve ipuçları
\end{itemize}

\section{BlackArch Linux Nedir?}
BlackArch, siber güvenlik araştırmacıları ve sızma testi uzmanları için geliştirilmekte olan bir GNU/Linux dağıtımıdır. \href{https://www.archlinux.org/}{ArchLinux} işletim sistemini taban olarak kabul etmiştir ve kullanıcılara BlackArch paketlerini kurma imkanı sağlamaktadır. Kullanıcılar bu paketleri grup olarak yükleyebileceği gibi, doğrudan da yükleyebilir.

Araçların yer aldığı paket deposu, Arch Linux işletim sisteminde olduğu gibi
\href{https://wiki.archlinux.org/index.php/Unofficial\_User\_Repositories}
{resmi olmayan kullanıcı deposu} olarak kullanılabilir. Doğal olarak hali hazırda kurulu olan Arch Linux işletim sistemine kurulum yapılabilir. Paket deposu eklemesi yapıldıktan sonra araçlar doğrudan veya gruplar halinde kurulabilir.

Şu anda depo içerisinde \href{https://www.blackarch.org/tools.html}{1300} araç bulunmaktadır. Depo kalitesinin korunması için tüm araçlar eklenmeden önce test edilmektedir.

% should quickly describe the testing methods/code review procedures etc

\section{BlackArch Linux Rehberi}
Yakında...

\section{Supported platforms}
Yakında...

\section{Dahil Olun}
BlackArch ekibi ile aşağıdaki yöntemlerden istediğiniz birisi ile iletişime geçebilirsiniz:

Web sitesi: \url{https://www.blackarch.org/}

E-Posta adresi: \href{mailto:blackarchlinux@gmail.com}{blackarchlinux@gmail.com}

IRC: \url{irc://irc.freenode.net/blackarch}

Twitter: \url{https://twitter.com/blackarchlinux}

Github: \url{https://github.com/Blackarch/}

%------------------%
%  Chapter 2       %
%------------------%


\chapter{Kullanıcı Rehberi}

\section{Kurulum}
Bu kısım altında BlackArch deposunun nasıl sisteme dahil edileceği ve paketlerin nasıl kurulacağı yer almaktadır. BlackArch depolarda yer alan derlenmiş paketleri desteklediği gibi kaynak kod ile paket kurulumunu da desteklemektedir.

BlackArch normal bir Arch Linux kurulumu ile uyumludur. Resmi olmayan kullanıcı deposu olarak düşünülebilir. Paket deposu yerine doğrudan ISO ile kurulum yapmak isterseniz \href{https://www.blackarch.org/downloads.html#iso}{Live ISO} kısmına bakabilirsiniz.


\subsection{Var olan Arch Linux üzerine kurulum}
\href{https://blackarch.org/strap.sh}{strap.sh} betiğini root yetkileri ile çalıştırınız ve aşağıdaki adımları takip ediniz. Hash değeri kontrolünü yapmayı unutmayınız.

\begin{lstlisting}
   curl -O https://blackarch.org/strap.sh
   sha1sum strap.sh # bu degere esit olmali: 86eb4efb68918dbfdd1e22862a48fda20a8145ff
   sudo ./strap.sh
\end{lstlisting}

Şimdi, ana paket listesinin güncellenmesi ve paketlerinizin senkron olması için aşağıdaki komutu kullanınız:
\begin{lstlisting}
  sudo pacman -Syyu
\end{lstlisting}


\subsection{Paket kurulumu}
Araçları blackarch deposundan kurabilirsiniz.
\begin{enumerate}
\item Kurulabilir tüm araçları listelemek için:
\begin{lstlisting}
  pacman -Sgg | grep blackarch | cut -d' ' -f2 | sort -u
\end{lstlisting}

\item Tüm araçları kurmak için:
\begin{lstlisting}
  pacman -S blackarch
\end{lstlisting}

\item Belirli bir kategoride yer alan araçları kurmak için:
\begin{lstlisting}
  pacman -S blackarch-<category>
\end{lstlisting}

\item Tüm blackarch kategorilerini görmek için:
\begin{lstlisting}
  pacman -Sg | grep blackarch
\end{lstlisting}

\end{enumerate}

\subsection{Kaynak koddan paket kurulumu}
Kurulumlara alternatif olarak BlackArch paketlerini doğrudan kaynak kod yardımı ile kurabilirsiniz. Bunun için PKGBUILD dosyalarına \href{https://github.com/BlackArch/blackarch/tree/master/packages}{paketler} dizini altından erişebilirsiniz. Tüm depoyu derlemek için \href{https://github.com/BlackArch/blackman}{Blackman} aracını kullanabilirsiniz.

\begin{itemize}
\item BlackArch paket deposu sisteminizde ekliyse doğrudan aşağıdaki komut yardımıyla Blackman aracını kurabilirsiniz.
\begin{lstlisting}
  pacman -S blackman
\end{lstlisting}

\item Blackman'i kaynak koddan kurmak için aşağıdaki adımları takip edebilirsiniz.
\begin{lstlisting}
  mkdir blackman
  cd blackman
  wget https://raw2.github.com/BlackArch/blackarch/master/packages/blackman/PKGBUILD
  # PKGBUILD dosyasinda herhangi bir sikinti olmadigini kontrol etmeyi unutmayiniz!
  makepkg -s
\end{lstlisting}

\item Blackman'i AUR aracılığı ile de kurabilirsiniz.
\begin{lstlisting}
  <AUR yardimci araciniz> -S blackman
\end{lstlisting}

\end{itemize}

\subsection{Temel BlackMan kullanımı}
Blackman her ne kadar alıştığınız diğer paket yöneticilerinin parametrelerinden farklı bir yapı kullansa da, gayet kolay bir kullanıma sahiptir. En temel kullanımı aşağıda gösterilmiştir.

\begin{itemize}
\item Paketi indir, derle ve kur:
\begin{lstlisting}
  sudo blackman -i paket
\end{lstlisting}

\item Paket grubunu indir, derle ve içerisinde yer alan tüm paketleri kur:
\begin{lstlisting}
  sudo blackman -g grup
\end{lstlisting}

\item Tüm BlackArch araçlarını indir, derle ve kur:
\begin{lstlisting}
  sudo blackman -a
\end{lstlisting}

\item BlackArch kategorilerini listeleme:
\begin{lstlisting}
  blackman -l
\end{lstlisting}

\item Kategori içerisinde yer alan araçları listeleme:
\begin{lstlisting}
  blackman -p category
\end{lstlisting}

\end{itemize}

\subsection{Doğrudan çalıştırılabilir ISO ile kurulum yapma}
BlackArch'ı doğrudan çalıştırılabilir ISO aracılığı ile de kurabilirsiniz. Bunun için \url{https://www.blackarch.org/download.html#iso} sayfasından bulabileceğiniz imaj ile aşağıdaki adımları takip edebilirsiniz. Adımları takip edebilmek için bilgisayarınızı indirdiğiniz imaj ile boot etmeniz gerekmektedir.

\begin{itemize}
\item blackarch-install-scripts paketi kurulumu:
\begin{lstlisting}
  sudo pacman -S blackarch-install-scripts
\end{lstlisting}

\item Çalıştırma
\begin{lstlisting}
  sudo blackarch-install
\end{lstlisting}

\end{itemize}

%------------------%
%  Chapter 3       %
%------------------%

\chapter{Geliştirici Rehberi}

\section{Arch Linux Paket Derleme ve Paket Depoları}

PKGBUILD dosyaları derleme betikleridir. Bu betikler makepkg(1) uygulamasına paketin nasıl oluşturulacağı bilgisini sağlamaktadır. PKGBUILD dosyaları Bash betikleri olarak yazılmaktadır.

Daha fazla bilgi için aşağıdaki kaynaklara bakabilirsiniz:

\begin{itemize}
\item \href{https://wiki.archlinux.org/index.php/Creating_Packages}{Arch Wiki: Paket Oluşturulması}
\item \href{https://wiki.archlinux.org/index.php/Makepkg}{Arch Wiki: makepkg}
\item \href{https://wiki.archlinux.org/index.php/PKGBUILD}{Arch Wiki: PKGBUILD}
\item \href{https://wiki.archlinux.org/index.php/Arch_Packaging_Standards}{Arch Wiki: Arch Paketleri Standartları}
\end{itemize}

\section{Blackarch PKGBUILD standartları}
AUR deposundaki paketler ile benzerliği korumak için PKGBUILD dosyaları ufak farklar haricinde benzer şekilde oluşturulmaktadır. Her bir paket en az blackarch grubuna dahil olmalıdır. Benzer şekilde bir paket birden fazla gruba dahil olabileceği gibi bir grup içerisinde de birden fazla paket bulunabilir.

\subsection{Gruplar}
Kullanıcılar kolay ve hızlı bir şekilde benzer paketleri kurabilmesi için çeşitli gruplamalar yapılmıştır. Gruplamanın getirdiği avantaj sayesinde "pacman -S <grup\_adi>" şeklinde bir komut ile tüm paketler kurulabilir.

\subsubsection{blackarch}
"blackarch" grubu diğer tüm grupların dahil olmak zorunda olduğu temel gruptur. Bu, tüm kullanıcıların tüm paketleri rahatlıkla kurmasını sağlar.

Bu grupta ne bulunabilir? Herşey.

\subsubsection{blackarch-anti-forensic}
Şifreleme, steganografi ve dosya özelliklerini düzenleme işlemlerini içeren adli bilişim işlemleri için kullanılan paketlerdir. 
İçerdiği tüm araçların amacı sistemde bilgi gizlemek için değişiklik yapmaktır.

Örnek: luks, TrueCrypt, Timestomp, dd, ropeadope, secure-delete

\subsubsection{blackarch-automation}
Araçların veya iş akışının düzenlenmesi için kullanılan paketlerdir.

Örnek: blueranger, tiger, wiffy

\subsubsection{blackarch-backdoor}
Zafiyetli sistemlerdeki açık backdoorlar veya exploitlerdir.

Örnek: backdoor-factory, rrs, weevely

\subsubsection{blackarch-binary}
Bazı formlardaki binary dosya işlemleri için kullanılan paketlerdir.

Örnek: binwally, packerid

\subsubsection{blackarch-bluetooth}
Bluetooth standardı(802.15.1) ile alakalı exploitlerdir.

Örnek: ubertooth, tbear, redfang

\subsubsection{blackarch-code-audit}
Zafiyet analizi için kaynak kod denetimi yapan paketlerdir.

Örnek: flawfinder, pscan

\subsubsection{blackarch-cracker}
Hash gibi kriptografik fonksiyonların cracklenmesinde kullanılan paketlerdir.

Örnek: hashcat, john, crunch

\subsubsection{blackarch-crypto}
Crack işlemleri haricindeki kriptografik işlemlerde kullanılan paketlerdir.

Örnek: ciphertest, xortool, sbd

\subsubsection{blackarch-database}
Herhangi bir seviyedeki veritabanı exploitlerini içeren palketlerdir.

Örnek: metacoretex, blindsql

\subsubsection{blackarch-debugger}
Belirli bir programın realtime'da yaptığı işi incelemeye yarayan paketlerdir.

Örnek: radare2, shellnoob

\subsubsection{blackarch-decompiler}
Derlenmiş programları kaynak koduna çevirmeye yarayan paketlerdir.

Örnek: flasm, jd-gui

\subsubsection{blackarch-defensive}
Kullanıcıyı zararlı yazılımlardan ve saldırılardan koruyan paketlerdir.

Örnek: arpon, chkrootkit, sniffjoke

\subsubsection{blackarch-disassembler}
Bu paketler blackarch-decompiler paketlerine benzer ve bir çok paket her iki grupta da bulunabilir. Bunların farkı ise binary dosyayı reverse ederek kaynak kodu yerine assembly çıktısı verirler.

Örnek: inguma, radare2

\subsubsection{blackarch-dos}
DoS (Denial of Service) saldırıları için kullanılan paketlerdir.

Örnek: 42zip, nkiller2

\subsubsection{blackarch-drone}
Fiziksel olarak drone'ları yönetmek için kullanılan paketlerdir.

Örnek: meshdeck, skyjack

\subsubsection{blackarch-exploitation}
Diğer program ve servisleri exploit ederken avantaj sağlayan paketlerdir.

Örnek: armitage, metasploit, zarp

\subsubsection{blackarch-fingerprint}
Parmak izi biyometrik okuyucuların exploit edilmesinde kullanılan paketlerdir.

Örnek: dns-map, p0f, httprint

\subsubsection{blackarch-firmware}
Firmware'deki zafiyetlerin exploit edilmesinde kullanılan paketlerdir.

Örnek: Henüz yok, en kısa sürede değiştirilecek.

\subsubsection{blackarch-forensic}
Fiziksel disk ve hafızadaki verilerin bulunmasında kullanılan paketlerdir.

Örnek: aesfix, nfex, wyd

\subsubsection{blackarch-fuzzer}
Fuzzing işlemlerinde kullanılan paketlerdir. Örneğin bu paketlerle hedefe random değerler göndererek sonucunda ne olduğunu inceleyebilirsiniz.

Örnek: msf, mdk3, wfuzz

\subsubsection{blackarch-hardware}
Fiziksel donanım ile yapılabilecek şeyleri yönetmek veya exploit etmek için kullanılan paketlerdir.

Örnek: arduino, smali

\subsubsection{blackarch-honeypot}
"Honeypot" görevi gören araçlardır. Honeypot'lar saldırganları tuzağa düşürmek için kullanılan programlardır.

Örnek: artillery, bluepot, wifi-honey

\subsubsection{blackarch-keylogger}
Bir sistemdeki klavye girdilerini saklayan paketlerdir.

Örnek: Henüz yok, en kısa sürede değiştirilecek.

\subsubsection{blackarch-malware}
Malware'lerin veya zararlı olabilecek yazılımların tespitinde kullanılan paketlerdir.

Örnek: malwaredetect, peepdf, yara

\subsubsection{blackarch-misc}
Başka herhangi bir gruba ait olmayan paketlerdir.

Örnek: oh-my-zsh-git, winexe, stompy

\subsubsection{blackarch-mobile}
Mobil platformların manipüle edilmesinde kullanılan paketlerdir.

Örnek: android-sdk-platform-tools, android-udev-rules

\subsubsection{blackarch-networking}
Network işlemlerinde kullanılan paketlerdir.

Örnek: Hemen hemen herşey

\subsubsection{blackarch-nfc}
NFC (near-field communications) işlemlerinde kullanılan paketlerdir.

Örnek: nfcutils

\subsubsection{blackarch-packer}
Packerları içeren paketlerdir.

/textif{Packerlar diğer yazılımlara zararlı yazılım gömebilen paketlerdir.}

Örnek: packerid

\subsubsection{blackarch-proxy}
Proxy görevi görevi gören uygulamalardır. Proxy'ler trafiği internetteki başka node'lar üzerinden geçiren yapılardır.

Örnek: burpsuite, ratproxy, sslnuke

\subsubsection{blackarch-recon}
Aktif zafiyet taraması yapan paketlerdir. Daha çok benzer paketlerin bir araya toplandığı şemsiye görevi görür.

Örnek: canri, dnsrecon, netmask

\subsubsection{blackarch-reversing}
Decompiler, disassembler gruplarının veya benzer programların bir araya toplandığı gruptur.

Örnek: capstone, radare2, zerowine

\subsubsection{blackarch-scanner}
Hedef sistemde zafiyet taraması yapan paketlerdir.

Örnek: scanssh, tiger, zmap

\subsubsection{blackarch-sniffer}
Ağ trafiğini analiz etmeye yarayan paketlerdir.

Örnek: hexinject, pytactle, xspy

\subsubsection{blackarch-social}
Öncelikli olarak sosyal ağ sitelerine saldırmaya yarayan paketlerdir.

Örnek: jigsaw, websploit

\subsubsection{blackarch-spoof}
Saldırganın kendisini gizleyerek kurbanı aldatma girişimlerinde kullandığı paketlerdir.

Examples: arpoison, lans, netcommander

\subsubsection{blackarch-threat-model}
Belirli bir senaryoda verilen tehdit modelini kaydetmek/raporlamak için kullanılan araçlardır.

Örnek: magictree

\subsubsection{blackarch-tunnel}
Verilen ağa tünnel ap trafiği oluşturmak için kullanılan paketlerdir.

Örnek: ctunnel, iodine, ptunnel

\subsubsection{blackarch-unpacker}
Çalıştırılabilir dosyalara yerleştirilmiş zararlı yazılımları çıkarmak için kullanılan paketlerdir.

Örnek: js-beautify

\subsubsection{blackarch-voip}
Voip programları ve protokolleri üzerinde işlem yapan paketlerdir.

Örnek: iaxflood, rtp-flood, teardown

\subsubsection{blackarch-webapp}
İnternet arayüzü olan uygulamaların işlemlerinde kullanılan paketlerdir.

Örnek: metoscan, whatweb, zaproxy

\subsubsection{blackarch-windows}
"wine" ile çalışan windows uygulamalarının bulunduğu paketlerdir.

Örnek: 3proxy-win32, pwdump, winexe

\subsubsection{blackarch-wireless}
Kablosuz ağ işlemleri yapan paketlerdir.

Örnek: airpwn, mdk3, wiffy

\section{Repository structure}
Ana BlackArch git deposuna buradan erişebilirsiniz:
\href{https://github.com/BlackArch/blackarch}{https://github.com/BlackArch/blackarch}.
Bir kaç tane ikincil depoya da erişmek için aşağıdaki linki kullanabilirsiniz:
\href{https://github.com/BlackArch}{https://github.com/BlackArch}.

Ana git deposunda üç önemli klasör bulunmaktadır:

\begin{itemize}
\item docs - Dökümantasyon.
\item packages - PKGBUILD dosyaları.
\item scripts - Basit kullanışlı scriptler.
\end{itemize}

\subsection{Betikler}
Aşağıda \verb|scripts/| klasöründe bulunan betikler için referanslar listelenmiştir:

\begin{itemize}
\item baaur - Yakında. Paketleri AUR'a yüklemek için kullanılacak.
\item babuild - Paketleri derler.
\item bachroot - Test için chroot'u yönetir.
\item baclean - Eski .pkg.tar.xz dosyalarını paket deposundan temizler.
\item baconflict - Yakında. \verb|scripts/conflicts|'leri düzeltecek.
\item bad-files - Derlenmiş paketlerdeki hatalı dosyaları bulur.
\item balock - Paket deposunu kilitler/açar.
\item banotify - Paket değişiklerinden IRC'yi haberdar eder.
\item barelease - Paketleri paket deposuna gönderir.
\item baright - BlackArch copyright bilgisini bastırır.
\item basign - Paketleri imzalar.
\item basign-key - Anahtarları imzalar.
\item blackman - Bir nevi pacman gibi davranır, fakat paketleri git'ten çeker.
(nrz'nin Blackman'inden farklıdır.).
\item check-groups - Grupları kontrol eder.
\item checkpkgs - Paketlerdeki hataları kontrol eder.
\item conflicts - Dosya çakışmalarını kontrol eder.
\item dbmod - Paket veritabanını düzenler.
\item depth-list - Bağımlılık(dependency) derinliğine göre liste oluşturur.
\item deptree - Yalnızca blackarch tarafından sunulan paketleri listeleyerek
    bağımlılık ağacını oluşturur.
\item get-blackarch-deps - Bir paket için blackarch bağımlılıklarını listeler.
\item get-official - Yayınlanacak resmi paketleri listeler.
\item list-loose-packages - Herhangi bir grupta yer almayan veya diğer paketlerle
    bağımlılığı olmayan paketleri listeler.
\item list-needed - Eksik bağımlılıkları listeler.
\item list-removed - Paket deposunda bulunan, git'te bulunmayan paketleri listeler.
\item list-tools - Araçları listeler.
\item outdated - Paket deposunda bulunan paketlerin versiyonlarının git deposuna göre
    eski olanları bulur.
\item pkgmod - Derlenmiş bir pakette değişiklik yapar.
\item pkgrel - Bir paketteki pkgrel'i artırır.
\item prep - PKGBUILD dosyalarını temizler ve hataları bulur.
\item sitesync - Bir paketin paket deposundaki local kopyası ile uzak sunucudaki
    kopyası arasında senkronizasyon sağlar.
\item size-hunt - Büyük paketleri arar.
\item source-backup - Paketlerin kaynak dosyalarını yedekler.
\end{itemize}

\section{Depoya katkıda bulunmak}
Bu bölüm BlackArch Linux projesine nasıl katkıda bulunabileceğinizi anlatmaktadır.
Küçük harf hatalası düzeltmelerinden yeni paketlere kadar her türlü pull request'ler
kabul edilmektedir. \\Yardımcı olmak, öneride bulunmak veya soru sormak
için bizimle iletişime geçebilirsiniz.
\\\\
Herkes katkıda bulunabilir. Tüm destekler değerlendirilecektir.

\subsection{Yardımcı kaynaklar}
Lütfen katkıda bulunmadan önce aşağıdaki kısmı okuyunuz:
\begin{itemize}
\item
\href{https://wiki.archlinux.org/index.php/Arch\_Packaging\_Standards)}{Arch
Paketleme Standartları}
\item \href{https://wiki.archlinux.org/index.php/Creating\_Packages}{Paket
Oluşturma}
\item \href{https://wiki.archlinux.org/index.php/PKGBUILD}{PKGBUILD}
\item \href{https://wiki.archlinux.org/index.php/Makepkg}{Makepkg}
\end{itemize}

\subsection{Katkı aşamaları}
BlackArch Linux projesine değişikliklerinizi göndermek için aşağıdaki
adımları takip edebilirsiniz:
\begin{enumerate}
\item Depoyu
\url{https://github.com/BlackArchLinux/blackarchlinux}
adresinden forklayın.
\item Gerekli dosyaları düzenleyin, (e.g. PKGBUILD, .patch dosyaları, vb).
\item Değişikliklerinizi commitleyin.
\item Değişikliklerinizi pushlayın.
\item Tercihen pull request ile değişikliklerinizi birleştirmemizi isteyin.
\end{enumerate}

\subsection{Örnek}
Aşağıdaki örnek yeni bir paketi BlackArch projesine göndermeyi göstermektedir.
\textbf{nfsshell} için önceden var olan PKGBUILD dosyasını almak için
\href{https://wiki.archlinux.org/index.php/yaourt}{yaourt} kullanıyoruz (isterseniz
pacaur'da kullanabilirsiniz.) ve ihtiyaçlarımıza göre ayarlıyoruz.

\subsubsection{PKGBUILD elde etme}
Yaourt veya pacaur kullanarak \textit{PKGBUILD} dosyasını alıyoruz:
\begin{lstlisting}
  user@blackarchlinux $ yaourt -G nfsshell
  ==> Download nfsshell sources
  x LICENSE
  x PKGBUILD
  x gcc.patch
  user@blackarchlinux $ cd nfsshell/
\end{lstlisting}

\subsubsection{PKGBUILD temizleme}
\textit{PKGBUILD} dosyasını temizleyerek biraz zaman kazanıyoruz:
\begin{lstlisting}
  user@blackarchlinux nfsshell $ ./blarckarch/scripts/prep PKGBUILD
  cleaning 'PKGBUILD'...
  expanding tabs...
  removing vim modeline...
  removing id comment...
  removing contributor and maintainer comments...
  squeezing extra blank lines...
  removing '|| return'...
  removing leading blank line...
  removing $pkgname...
  removing trailing whitespace...
\end{lstlisting}

\subsubsection{PKGBUILD ayarlama}
\textit{PKGBUILD} dosyasını ayarlıyoruz:
\begin{lstlisting}
  user@blackarchlinux nfsshell $ vi PKGBUILD
\end{lstlisting}

\subsubsection{Paket derleme}
Paketi derliyoruz:
\begin{lstlisting}user@blackarchlinux nfsshell $ makepkg -sf
==> Making package: nfsshell 19980519-1 (Mon Dec  2 17:23:51 CET 2013)
==> Checking runtime dependencies...
==> Checking buildtime dependencies...
==> Retrieving sources...
-> Downloading nfsshell.tar.gz...
% Total    % Received % Xferd  Average Speed   Time    Time     Time
CurrentDload  Upload   Total   Spent    Left  Speed100 29213  100 29213    0
0  48150      0 --:--:-- --:--:-- --:--:-- 48206
-> Found gcc.patch
-> Found LICENSE
...
<lots of build process and compiler output here>
...
==> Leaving fakeroot environment.
==> Finished making: nfsshell 19980519-1 (Mon Dec  2 17:23:53 CET 2013)
\end{lstlisting}

\subsubsection{Paketin kurulumu ve test edilmesi}
Paketi kurup test ediyoruz:
Install and test the package:
\begin{lstlisting}
  user@blackarchlinux nfsshell $ pacman -U nfsshell-19980519-1-x86_64.pkg.tar.xz
  user@blackarchlinux nfsshell $ nfsshell # test it
\end{lstlisting}

\subsubsection{Git'e gönderme}
Paketi commitleyip git'e gönderiyoruz:
\begin{lstlisting}user@blackarchlinux nfsshell $ cd /blackarchlinux/packages
user@blackarchlinux ~/blackarchlinux/packages $ mv ~/nfsshell .
user@blackarchlinux ~/blackarchlinux/packages $ git commit -am nfsshell && git push
\end{lstlisting}

\subsubsection{İstek gönderme}
\href{https://github.com/}{github.com} üzerinden pull request oluşturuyoruz:
\begin{lstlisting}
  firefox https://github.com/<contributor>/blackarchlinux
\end{lstlisting}

\subsubsection{Güncellemeleri takip etmek}
Eğer kendi forkladığınız depo üzerinde çalışıyorsanız ve ana ba deposunu uzak depo olarak belirlediyseniz
güncellemeleri takip etmek için yapılacaklar:
\begin{lstlisting}
  user@blackarchlinux ~/blackarchlinux $ git remote -v
  origin <the url of your fork> (fetch)
  origin <the url of your fork> (push)
  user@blackarchlinux ~/blackarchlinux $ git remote add upstream https://github.com/blackarch/blackarch
  user@blackarchlinux ~/blackarchlinux $ git remote -v
  origin <the url of your fork> (fetch)
  origin <the url of your fork> (push)
  upstream https://github.com/blackarch/blackarch (fetch)
  upstream https://github.com/blackarch/blackarch (push)
\end{lstlisting}

Git varsayılan olarak origin'e gönderir ama git config dosyanızın doğru ayarlandığından
emin olmanız gerekir. Normal commitler sırasında bu problem olmaz fakat eğer
uzak sunucuya göndermek isterseniz bu mümkün olmayacaktır.

Eğer yapabiliyorsanız \textit{git@github.com:blackarch/blackarch.git}
kullanarak commitlerinizi yapmanız daha başarılı olacaktır. Fakat bu
konuda tercih size kalmış.

\subsection{İstekler}
\begin{enumerate}
\item \textbf{Maintainer} ya da \textbf{Contributor} isimlerini yorum olarak \textit{PKGBUILD} 
dosyalarına eklemeyin. Maintainer ve contributor isimlerini BlackArch rehberindeki AUTHORS
bölümüne ekleyebilirsiniz.
\item Tutarlılığa uyması adına depoda bulunan başka bir \textit{PKGBUILD} dosyasının genel stilini kullanınız.
Ayrıca girintilerde iki boşluk kullanınız.
\end{enumerate}

\subsection{Genel ipuçları}
\href{http://wiki.archlinux.org/index.php/Namcap}{namcap} ile paket hatalarını kontrol edebilirsiniz.

%------------------%
%  Chapter 4       %
%------------------%

\chapter{Araç Rehberi}
Yakında...

\section{Yakında}
Yakında...

%%% APPENDIX %%%
\appendix
%%%%%%%%%%%%%%%%%%%%%%%%%%%%%%%%%%%%%%%%%%%%%%%%%%%%%%%%%%%%%%%%%%%%%%%%%%%%%%%%
%                                                                              %
% BlackArch Linux appendix template                                            %
%                                                                              %
%%%%%%%%%%%%%%%%%%%%%%%%%%%%%%%%%%%%%%%%%%%%%%%%%%%%%%%%%%%%%%%%%%%%%%%%%%%%%%%%

\appendix

\chapter{Appendix}

\section{FAQs}

\section{AUTHORS}
The following people have contributed directly to BlackArch:
\begin{itemize}
\item Tyler Bennnett (tylerb@trix2voip.com)
\item fnord0 (fnord0@riseup.net)
\item nrz (nrz@nullsecurity.net)
\item Ellis Kenyo (elken.tdos@gmail.com)
\item CaledoniaProject (the.warl0ck.1989@gmail.com)
\item sudokode (sudokode@gmail.com)
\item Valentin Churavy (v.churavy@gmail.com)
\item Boy Sandy Gladies Arriezona (reno.esper@gmail.com)
\item Mathias Nyman
\item Johannes Löthberg (demizide@gmail.com)
\item Thiago da Silva Teixeira (teixeira.zeus@gmail.com)
\end{itemize}
The following people contributed directly to ArchPwn, which has been merged
into BlackArch:
\begin{itemize}
\item Francesco Piccinno (stack.box@gmail.com)
\item jensp (jens@jenux.homelinux.org)
\item Valentin Churavy (v.churavy@gmail.com)
\end{itemize}
We have taken build code from the following people:
\begin{itemize}
\item  3ED (krzysztof1987@gmail.com)
\item  AUR Perl (aurperl@juster.info)
\item  Aaron Griffin (aaron@archlinux.org)
\item  Abakus (java5@arcor.de)
\item  Adam Wolk (netprobe@gmail.com)
\item  Aleix Pol (aleixpol@kde.org)
\item  Aleshus (aleshusi@gmail.com)
\item  Alessandro Pazzaglia (jackdroido@gmail.com)
\item  Alessandro Sagratini (ale\_sagra@hotmail.com)
\item  Alex Cartwright (alexc223@googlemail.com)
\item  Alexander De Sousa (archaur.xandy21@spamgourmet.com)
\item  Alexander Rødseth (rodseth@gmail.com)
\item  Allan McRae (allan@archlinux.org)
\item  AmaN (gabroo.punjab.da@gmail.com)
\item  Andre Klitzing (aklitzing@online.de)
\item  Andrea Scarpino (andrea@archlinux.org)
\item  Andreas Schönfelder (passtschu@freenet.de)
\item  Andrej Gelenberg (andrej.gelenberg@udo.edu)
\item  Angel Velasquez (angvp@archlinux.org)
\item  Antoine Lubineau (antoine@lubignon.info)
\item  Anton Bazhenov (anton.bazhenov@gmail.com)
\item  Arkham (arkham@archlinux.us)
\item  Arthur Danskin (arthurdanskin@gmail.com)
\item  Balda (balda@balda.ch)
\item  Balló György (ballogyor+arch@gmail.com)
\item  Bartek Piotrowski (barthalion@gmail.com)
\item  Bartosz Feński (fenio@debian.org)
\item  Bartłomiej Piotrowski (nospam@bpiotrowski.pl)
\item  Bogdan Szczurek (thebodzio@gmail.com)
\item  Brad Fanella (bradfanella@archlinux.us)
\item  Brian Bidulock (bidulock@openss7.org)
\item  C Anthony Risinger (anthony@xtfx.me)
\item  CRT (crt.011@gmail.com)
\item  Can Celasun (dcelasun@gmail.com)
\item  Chaniyth (chaniyth@yahoo.com)
\item  Chris Brannon (cmbrannon79@gmail.com)
\item  Chris Giles (Chris.G.27@gmail.com) \& daschu117
\item  Christoph Siegenthaler (csi@gmx.ch)
\item  Christoph Zeiler (archNOSPAM@moonblade.org)
\item  Clément DEMOULINS (clement@archivel.fr)
\item  Corrado Primier (bardo@aur.archlinux.org)
\item  Daenyth (Daenyth+Arch@gmail.com)
\item  Dale Blount (dale@archlinux.org)
\item  Damir Perisa (damir.perisa@bluewin.ch)
\item  Dan Fuhry (dan@fuhry.us)
\item  Dan Serban (dserban01@yahoo.com)
\item  Daniel A. Campoverde Carrión
\item  Daniel Golle
\item  Daniel Griffiths (ghost1227@archlinux.us)
\item  Daniel J Griffiths (ghost1227@archlinux.us)
\item  Daniel Micay (danielmicay@gmail.com)
\item  Dave Reisner (dreisner@archlinux.org)
\item  Dawid Wrobel (cromo@klej.net)
\item  Devaev Maxim (mdevaev@gmail.com)
\item  Devin Cofer (ranguvar@archlinux.us)
\item  DigitalPathogen (aur@InfoSecResearchLabs.co.uk)
\item  DigitalPathogen (aur@digitalpathogen.co.uk)
\item  Dmitry A. Ilyashevich (dmitry.ilyashevich@gmail.com)
\item  Dominik Heidler (dheidler@gmail.com)
\item  DrZaius (lou@fakeoutdoorsman.com)
\item  Ebubekir KARUL (ebubekirkarul@yandex.com)
\item  Eduard "bekks" Warkentin (eduard.warkentin@gmail.com)
\item  Elmo Todurov (todurov@gmail.com)
\item  Emmanuel Gil Peyrot (linkmauve@linkmauve.fr)
\item  Eric Belanger (eric@archlinux.org)
\item  Ermak (ermak@email.it)
\item  Evangelos Foutras (evangelos@foutrelis.com)
\item  Fabian Melters (melters@gmail.com)
\item  Fabiano Furtado (fusca14@gmail.com)
\item  Federico Quagliata (ntp@quaqo.org)
\item  Firmicus (francois.archlinux@org)
\item  Florian Pritz (bluewind@jabber.ccc.de)
\item  Florian Pritz (flo@xinu.at)
\item  Francesco Piccinno (stack.box@gmail.com)
\item  François Charette (francois@archlinux.org)
\item  Gaetan Bisson (bisson@archlinux.org)
\item  Geoffroy Carrier (geoffroy.carrier@koon.fr)
\item  Georg Grabler (STiAT)
\item  George Hilliard (gh403@msstate.edu)
\item  Gerardo Exequiel Pozzi (vmlinuz386@yahoo.com.ar)
\item  Gilles CHAUVIN (gcnweb@gmail.com)
\item  Giovanni Scafora (giovanni@archlinux.org)
\item  Gordin (9ordin@gmail.com)
\item  Guillaume ALAUX (guillaume@archlinux.org)
\item  Guillermo Ramos (0xwille@gmail.com)
\item  Gustavo Alvarez (sl1pkn07@gmail.com)
\item  Hugo Doria (hugo@archlinux.org)
\item  Hyacinthe Cartiaux (hyacinthe.cartiaux@free.fr)
\item  James Fryman (jfryman@gmail.com)
\item  Jan "heftig" Steffens (jan.steffens@gmail.com)
\item  Jan de Groot (jgc@archlinux.org)
\item  Jaroslav Lichtblau (dragonlord@aur.archlinux.org)
\item  Jaroslaw Swierczynski (swiergot@aur.archlinux.org)
\item  Jason Chu (jason@archlinux.org)
\item  Jason R Begley (jayray@digitalgoat.com)
\item  Jason Rodriguez
\item  Jason St. John (jstjohn@purdue.edu)
\item  Jawmare (victor2008@gmail.com)
\item  Jeff Mickey (jeff@archlinux.org)
\item  Jens Pranaitis (jens@chaox.net)
\item  Jens Pranaitis (jens@jenux.homelinux.org)
\item  Jinx (jinxware@gmail.com)
\item  John D Jones III (jnbek1972@gmail.com)
\item  John Proctor (jproctor@prium.net)
\item  Jon Bergli Heier (snakebite@jvnv.net)
\item  Jonas Heinrich
\item  Jonathan Steel (jsteel@aur.archlinux.org)
\item  Joris Steyn (jorissteyn@gmail.com)
\item  Josh VanderLinden (arch@cloudlery.com)
\item  Jozef Riha (jose1711@gmail.com)
\item  Judd Vinet (jvinet@zeroflux.org)
\item  Juergen Hoetzel (jason@archlinux.org)
\item  Juergen Hoetzel (juergen@archlinux.org)
\item  Justin Davis (jrcd83@gmail.com)
\item  Kaiting Chen (kaitocracy@gmail.com)
\item  Kaos
\item  Kevin Piche (kevin@archlinux.org)
\item  Kory Woods (kory@virlo.net)
\item  Kyle Keen (keenerd@gmail.com)
\item  Larry Hajali (larryhaja@gmail.com)
\item  LeCrayonVert
\item  Le\_suisse (lesuisse.dev+aur@gmail.com)
\item  Lekensteyn (lekensteyn@gmail.com)
\item  Limao Luo (luolimao+AUR@gmail.com)
\item  Lucien Immink
\item  Lukas Fleischer (archlinux@cryptocrack.de)
\item  Manolis Tzanidakis
\item  Marcin "avalan" Falkiewicz (avalatron@gmail.com)
\item  Mariano Verdu (verdumariano@gmail.com)
\item  Marti Raudsepp (marti@juffo.org)
\item  MatToufoutu (mattoufootu@gmail.com)
\item  Matthew Sharpe (matt.sharpe@gmail.com)
\item  Mauro Andreolini (mauro.andreolini@unimore.it)
\item  Max Pray a.k.a. Synthead (synthead@gmail.com)
\item  Max Roder (maxroder@web.de)
\item  Maxwell Pray a.k.a. Synthead (synthead@gmail.com)
\item  Maxwell Pray a.k.a. Synthead (synthead1@gmail.com)
\item  Mech (tiago.bmp@gmail.com)
\item  Michael Düll (mail@akurei.me)
\item  Michael P (ptchinster@archlinux.us)
\item  Michal Krenek (mikos@sg1.cz)
\item  Michal Zalewski (lcamtuf@coredump.cx)
\item  Miguel Paolino (mpaolino@gmail.com)
\item  Miguel Revilla (yo@miguelrevilla.com)
\item  Mike Roberts (noodlesgc@gmail.com)
\item  Mike Sampson (mike@sambodata.com)
\item  Nassim Kacha (nassim.kacha@gmail.com)
\item  Nicolas Pouillard (nicolas.pouillard@gmail.com)
\item  Nicolas Pouillard https://nicolaspouillard.fr
\item  Niklas Schmuecker
\item  Oleander Reis (oleander@oleander.cc)
\item  Olivier Le Moal (mail@olivierlemoal.fr)
\item  Olivier Médoc "oliv" (o\_medoc@yahoo.fr)
\item  Pascal E. (archlinux@hardfalcon.net)
\item  Patrick Leslie Polzer (leslie.polzer@gmx.net)
\item  Paul Mattal (paul@archlinux.org)
\item  Paul Mattal (pjmattal@elys.com)
\item  Pengyu CHEN (cpy.prefers.you@gmail.com)
\item  Peter Wu (lekensteyn@gmail.com)
\item  Philipp 'TamCore' B. (philipp@tamcore.eu)
\item  Pierre Schmitz (pierre@archlinux.de)
\item  Pranay Kanwar (pranay.kanwar@gmail.com)
\item  Pranay Kanwar (warl0ck@metaeye.org)
\item  PyroPeter (googlemail@com.abi1789)
\item  PyroPeter (googlemail.com@abi1789)
\item  Ray Rashif (schiv@archlinux.org)
\item  Remi Gacogne
\item  Renan Fernandes (renan@kauamanga.com)
\item  Richard Murri (admin@richardmurri.com)
\item  Roberto Alsina (ralsina@kde.org)
\item  Robson Peixoto (robsonpeixoto@gmail.com)
\item  Roel Blaauwgeers (roel@ttys0.nl)
\item  Rorschach (r0rschach@lavabit.com)
\item  Ruben Schuller (shiml@orgizm.net)
\item  Rudy Matela (rudy@matela.com)
\item  Ryon Sherman (ryon.sherman@gmail.com)
\item  Sabart Otto \item  Seberm (seberm@gmail.com)
\item  SakalisC (chrissakalis@gmail.com)
\item  Sam Stuewe (halosghost@archlinux.info)
\item  SanskritFritz (SanskritFritz@gmail.com)
\item  Sarah Hay (sarahhay@mb.sympatico)
\item  Sebastian Benvenuti (sebastianbenvenuti@gmail.com)
\item  Sebastian Nowicki (sebnow@gmail.com)
\item  Sebastien Duquette (ekse.0x@gmail.com)
\item  Sebastien LEDUC (sebastien@sleduc.fr)
\item  Sebastien Leduc (sebastien@sleduc.fr)
\item  Sergej Pupykin (pupykin.s+arch@gmail.com)
\item  Sergio Rubio (rubiojr@biondofu.net)
\item  Sheng Yu (magicfish1990@gmail.com)
\item  Simon Busch (morphis@gravedo.de)
\item  Simon Legner (Simon.Legner@gmail.com)
\item  Sirat18 (aur@sirat18.de)
\item  SpepS (dreamspepser@yahoo.it)
\item  Spider.007 (archPackage@spider007.net)
\item  Stefan Seering
\item  Stephane Travostino (stephane.travostino@gmail.com)
\item  Stéphane Gaudreault (stephane@archlinux.org)
\item  Sven Kauber (celeon@gmail.com)
\item  Sven Schulz (omee@archlinux.de)
\item  Sébastien Duquette (ekse.0x@gmail.com)
\item  Sébastien Luttringer (seblu@archlinux.org)
\item  TDY (tdy@gmx.com)
\item  Teemu Rytilahti (tpr@iki.fi)
\item  Testuser\_01
\item  Thanx (thanxm@gmail.com)
\item  Thayer Williams (thayer@archlinux.org)
\item  Thomas S Hatch (thatch45@gmail.com)
\item  Thorsten Töpper
\item  Tilmann Becker (tilmann.becker@freenet.de)
\item  Timothy Redaelli (timothy.redaelli@gmail.com)
\item  Timothée Ravier (tim@siosm.fr)
\item  Tino Reichardt
\item  Tobias Kieslich (tobias@justdreams.de)
\item  Tobias Powalowski (tpowa@archlinux.org)
\item  Tom K (tomk@runbox.com)
\item  Tom Newsom (Jeepster@gmx.co.uk)
\item  Tomas Lindquist Olsen (tomas@famolsen.dk)
\item  Travis Willard (travisw@wmpub.ca)
\item  Valentin Churavy (v.churavy@gmail.com)
\item  ViNS (gladiator@fastwebnet.it)
\item  Vlatko Kosturjak (kost@linux.hr)
\item  Wes Brown (wesbrown18@gmail.com)
\item  William Rea (sillywilly@gmail.com)
\item  Xavier Devlamynck (magicrhesus@ouranos.be)
\item  Xiao\item Long Chen (chenxiaolong@cxl.epac.to)
\item  aeolist (aeolist@hotmail.com)
\item  ality@pbrane.org
\item  astaroth (astaroth\_@web.de)
\item  bender02@archlinux.us
\item  billycongo (billycongo@gmail.com)
\item  bslackr (brendan@vastactive.com)
\item  cbreaker (cbreaker@tlen.pl)
\item  chimeracoder (dev@chimeracoder.net)
\item  damir (damir@archlinux.org)
\item  danitool
\item  darkapex (me@jailuthra.in)
\item  daronin
\item  dkaylor (dpkaylor@gmail.com)
\item  dobo (dobo90\_at\_gmail@com)
\item  dorphell (dorphell@archlinux.org)
\item  evr (evanroman@at.gmail)
\item  fnord0 (fnord0@riseup.net)
\item  fxbru (frxbru@gmail)
\item  hcar
\item  icarus (icarus.roaming@gmail.com)
\item  iceman (icemanf@gmail.com)
\item  kastor (kastor@fobos.org)
\item  kfgz (kfgz@interia.pl)
\item  linuxSEAT (linuxSEAT@gmail.com)
\item  m4xm4n (max@maxfierke.com)
\item  mar77i (mysatyre@gmail.com)
\item  marc0s (marc0s@fsfe.org)
\item  mickael9 (mickael9@gmail.com)
\item  nblock (nblock@archlinux.us)
\item  nofxx (x@nofxx.com)
\item  onny (onny@project
\item  pootzko (pootzko@gmail.com)
\item  revel (revel@muub.net)
\item  rich\_o (rich\_o@lavabit.com)
\item  s1gma (s1gma@mindslicer.com)
\item  sandman (r.coded@gmail.com)
\item  sebikul (sebikul@gmail.com)
\item  sh0 (mee@sh0.org)
\item  shild (sxp@bk.ru)
\item  simo (simo@archlinux.org)
\item  snuo
\item  sudokode (sudokode@gmail.com)
\item  tobias (tobias@archlinux.org)
\item  trashstar (trash@ps3zone.org)
\item  unexist (unexist@subforge.org)
\item  untitled (rnd0x00@gmail.com)
\item  virtuemood (virtue@archlinux.us)
\item  wido (widomaker2k7@gmail.com)
\item  wodim (neikokz@gmail.com)
\item  yannsen (ynnsen@gmail.com)
\end{itemize}


\end{document}

%%% EOF %%%
